% GENERATED FILE. Edit canonical lessons, then run npm run generate:books.
% canonical-source-hash: fnv1a64:fb4abe0b1449bd66
% canonical-lessons: GE-C528-unordnung, GE-C528-ausmisten, GE-C528-abfall, GE-C528-muelltonne, GE-C528-entsorgen

\chapter{Müll und Unordnung}
\label{ch:ge-mull-und-unordnung}
\hlchaptermodality{528}

\begin{chapteropening}
\textbf{By the end of this chapter:} \emph{I can talk about rubbish and clearing out.}

Complete the last lesson of chapter 528: I can talk about rubbish and clearing out.
\end{chapteropening}

\section[die Unordnung]{die Unordnung — mess, disorder}

\label{lesson:GE-C528-unordnung}



\begin{quote}
Before the new one: say the German for to iron, then the German for to fold.
\end{quote}

\subsection*{You'll want to know: die Unordnung}
\textbf{die Unordnung} — \textquotedblleft{}mess, disorder\textquotedblright{}.

\begin{quote}
In seinem Zimmer herrscht immer Unordnung.
\end{quote}

\begin{quote}
\emph{His room is forever in a mess.}
\end{quote}

No plural.

\begin{cousinweb}
\emph{Unordnung} is built on \emph{un-} \textquotedblleft{}not\textquotedblright{} and \emph{Ordnung} \textquotedblleft{}order\textquotedblright{}, inherited from Old High German \emph{ordinunga} (EtymDB), built on a verb borrowed from Latin \emph{ōrdināre}, from \emph{ōrdō} \textquotedblleft{}row, order\textquotedblright{}.
\end{cousinweb}

\begin{grammarlens}[title={What you've built}]
One more word for telling, explaining and describing at B1.
\end{grammarlens}

\subsection*{Your turn}
\begin{itemize}
\raggedright
  \item \emph{Say it:} \emph{die Unordnung}
  \item \emph{Say it:} \emph{die Unordnung}, once more
  \item \emph{Say it:} a sentence of your own with \emph{die Unordnung}
  \item \emph{Recall:} say the German for to iron, then the German for to fold, then say \emph{die Unordnung} again
\end{itemize}

\begin{tcolorbox}[breakable,colback=teal!4,colframe=teal!35!black,title={Before you move on}]
What does die Unordnung mean? (\textquotedblleft{}Mess, disorder\textquotedblright{}.) Say it once more.
\end{tcolorbox}
\section[ausmisten]{ausmisten — to clear out}

\label{lesson:GE-C528-ausmisten}



\begin{quote}
Before the new one: say the German for to fold, then the German for mess.
\end{quote}

\subsection*{You'll want to know: ausmisten}
\textbf{ausmisten} — \textquotedblleft{}to clear out\textquotedblright{}.

\begin{quote}
Am Wochenende haben wir den Keller ausgemistet.
\end{quote}

\begin{quote}
\emph{At the weekend we cleared out the cellar.}
\end{quote}

Separable: \emph{wir misten aus}. It first meant to muck out a stable.

\begin{cousinweb}
\emph{ausmisten} is built on \emph{aus} \textquotedblleft{}out\textquotedblright{} and \emph{misten} \textquotedblleft{}to muck out\textquotedblright{}, from \emph{Mist} \textquotedblleft{}manure\textquotedblright{}, inherited from Old High German \emph{mist} (EtymDB).
\end{cousinweb}

\begin{grammarlens}[title={What you've built}]
One more word for telling, explaining and describing at B1.
\end{grammarlens}

\subsection*{Your turn}
\begin{itemize}
\raggedright
  \item \emph{Say it:} \emph{ausmisten}
  \item \emph{Say it:} \emph{ausmisten}, once more
  \item \emph{Say it:} a sentence of your own with \emph{ausmisten}
  \item \emph{Recall:} say the German for to fold, then the German for mess, then say \emph{ausmisten} again
\end{itemize}

\begin{tcolorbox}[breakable,colback=teal!4,colframe=teal!35!black,title={Before you move on}]
What does ausmisten mean? (\textquotedblleft{}To clear out\textquotedblright{}.) Say it once more.
\end{tcolorbox}
\section[der Abfall]{der Abfall — rubbish, waste}

\label{lesson:GE-C528-abfall}



\begin{quote}
Before the new one: say the German for mess, then the German for to clear out.
\end{quote}

\subsection*{You'll want to know: der Abfall}
\textbf{der Abfall} — \textquotedblleft{}rubbish, waste\textquotedblright{}.

\begin{quote}
Im Park liegt viel Abfall auf dem Boden.
\end{quote}

\begin{quote}
\emph{There's a lot of rubbish on the ground in the park.}
\end{quote}

Plural: \emph{die Abfälle}. In speech also \emph{der Müll}.

\begin{cousinweb}
\emph{Abfall} is built on \emph{ab} \textquotedblleft{}off\textquotedblright{} and \emph{Fall} \textquotedblleft{}fall\textquotedblright{}, inherited from Old High German \emph{val} (EtymDB): what falls off.
\end{cousinweb}

\begin{grammarlens}[title={What you've built}]
One more word for telling, explaining and describing at B1.
\end{grammarlens}

\subsection*{Your turn}
\begin{itemize}
\raggedright
  \item \emph{Say it:} \emph{der Abfall}
  \item \emph{Say it:} \emph{der Abfall}, once more
  \item \emph{Say it:} a sentence of your own with \emph{der Abfall}
  \item \emph{Recall:} say the German for mess, then the German for to clear out, then say \emph{der Abfall} again
\end{itemize}

\begin{tcolorbox}[breakable,colback=teal!4,colframe=teal!35!black,title={Before you move on}]
What does der Abfall mean? (\textquotedblleft{}Rubbish, waste\textquotedblright{}.) Say it once more.
\end{tcolorbox}
\section[die Mülltonne]{die Mülltonne — dustbin, wheelie bin}

\label{lesson:GE-C528-muelltonne}



\begin{quote}
Before the new one: say the German for to clear out, then the German for rubbish.
\end{quote}

\subsection*{You'll want to know: die Mülltonne}
\textbf{die Mülltonne} — \textquotedblleft{}dustbin, wheelie bin\textquotedblright{}.

\begin{quote}
Am Dienstag stellen wir die Mülltonne an die Straße.
\end{quote}

\begin{quote}
\emph{On Tuesday we put the bin out on the street.}
\end{quote}

Plural: \emph{die Mülltonnen}.

\begin{cousinweb}
\emph{Mülltonne} is built on \emph{Müll} \textquotedblleft{}rubbish\textquotedblright{} and \emph{Tonne} \textquotedblleft{}barrel, drum\textquotedblright{}, inherited from Old High German \emph{tunna} (EtymDB).
\end{cousinweb}

\begin{grammarlens}[title={What you've built}]
One more word for telling, explaining and describing at B1.
\end{grammarlens}

\subsection*{Your turn}
\begin{itemize}
\raggedright
  \item \emph{Say it:} \emph{die Mülltonne}
  \item \emph{Say it:} \emph{die Mülltonne}, once more
  \item \emph{Say it:} a sentence of your own with \emph{die Mülltonne}
  \item \emph{Recall:} say the German for to clear out, then the German for rubbish, then say \emph{die Mülltonne} again
\end{itemize}

\begin{tcolorbox}[breakable,colback=teal!4,colframe=teal!35!black,title={Before you move on}]
What does die Mülltonne mean? (\textquotedblleft{}Dustbin, wheelie bin\textquotedblright{}.) Say it once more.
\end{tcolorbox}
\section[entsorgen]{entsorgen — to dispose of}

\label{lesson:GE-C528-entsorgen}



\begin{quote}
Before the new one: say the German for rubbish, then the German for dustbin.
\end{quote}

\subsection*{You'll want to know: entsorgen}
\textbf{entsorgen} — \textquotedblleft{}to dispose of\textquotedblright{}.

\begin{quote}
Alte Batterien muss man richtig entsorgen.
\end{quote}

\begin{quote}
\emph{Old batteries have to be disposed of properly.}
\end{quote}

Inseparable: \emph{hat entsorgt}.

\begin{cousinweb}
\emph{entsorgen} is built on \emph{ent-} \textquotedblleft{}away\textquotedblright{} and \emph{Sorge} \textquotedblleft{}care, worry\textquotedblright{}, inherited from Old High German \emph{sorga}, from Proto-Germanic \emph{surgō} (EtymDB): to take away a worry.
\end{cousinweb}

\begin{grammarlens}[title={What you've built}]
One more word for telling, explaining and describing at B1.
\end{grammarlens}

\subsection*{Your turn}
\begin{itemize}
\raggedright
  \item \emph{Say it:} \emph{entsorgen}
  \item \emph{Say it:} \emph{entsorgen}, once more
  \item \emph{Say it:} a sentence of your own with \emph{entsorgen}
  \item \emph{Recall:} say the German for rubbish, then the German for dustbin, then say \emph{entsorgen} again
\end{itemize}

\begin{tcolorbox}[breakable,colback=teal!4,colframe=teal!35!black,title={Before you move on}]
What does entsorgen mean? (\textquotedblleft{}To dispose of\textquotedblright{}.) Say it once more.
\end{tcolorbox}
